% Generated draft: requires mathematical and source review.
\documentclass[11pt]{article}
\usepackage[margin=25mm]{geometry}
\usepackage{fontspec}
\setmainfont{Latin Modern Roman}
\usepackage{amsmath,amssymb,mathtools,unicode-math}
\setmathfont{Latin Modern Math}
\usepackage{enumitem,needspace,xurl,hyperref}
\setlength{\parindent}{0pt}
\setlength{\parskip}{5pt}
\setlength{\emergencystretch}{3em}
\newcommand{\N}{\mathbb{N}}
\newcommand{\Z}{\mathbb{Z}}
\newcommand{\Q}{\mathbb{Q}}
\newcommand{\R}{\mathbb{R}}
\newcommand{\C}{\mathbb{C}}
\newcommand{\F}{\mathbb{F}}
\newcommand{\A}{\mathbb{A}}
\newcommand{\PP}{\mathbb{P}}
\newcommand{\E}{\mathbb{E}}
\newcommand{\eps}{\varepsilon}
\newcommand{\dd}{\,\mathrm{d}}
\begin{document}
\section*{UnsolvedMath ID: 1200019}
\textbf{Source problem code:} catalogue entry
\subsection{Definitions and mathematical statement}
Fix a degree $d$. To specify the standard finite-graph side, let $G_N$ be uniform among labelled simple $d$-regular graphs on $N$ vertices, along admissible sizes $N>d$ with $dN$ even. The source does not specify this model rather than a configuration multigraph model, nor the range of $d$.

For integers $r\geq0$ and $k\geq1$, and a coloring $c:V(G)\to\{1,\ldots,k\}$ of a finite graph, let $\mu_{G,c,r}$ be the probability distribution of the rooted colored radius-$r$ neighborhood of a uniformly chosen vertex. Put
\[
 Q_{r,k}(G)=\{\mu_{G,c,r}:c:V(G)\to\{1,\ldots,k\}\}.
\]
For fixed degree bound, $r,k$, these are subsets of a finite-dimensional probability simplex. Local-global convergence compares these sets for every $r,k$ in Hausdorff distance induced by total variation. Unlike local weak convergence, it records all finite colorings, not just uncolored neighborhoods.

The question asks whether uniformly random $d$-regular graphs converge in this sense to the weak closure of i.i.d. processes. The source does not identify the underlying process space, its symmetry/invariance convention, or whether ``i.i.d. processes'' means independent label processes themselves or their equivariant factors. For the usual candidate factor-of-i.i.d. setting, one starts with independent labels on the vertices of the infinite $d$-regular tree and applies a measurable automorphism-equivariant coloring map; weak closure then means closure of their laws under convergence of finite-neighborhood distributions. This describes a possible intended target, not a substitution for the source's unresolved terminology.

Once a target family $\mathcal P_{d,k}$ and a random-convergence mode have been fixed, the corresponding assertion would require, for every fixed $r,k$,
\[
 d_H\bigl(Q_{r,k}(G_N),\ \{\text{radius-}r\text{ marginal of }\nu:
                     \nu\in\mathcal P_{d,k}\}\bigr)\longrightarrow0
\]
as $N\to\infty$ through admissible sizes. The missing conventions prevent this conditional display from being identified as a uniquely determined formal version of the original question.
\subsection{Short English statement}
Does the complete collection of finite-coloring statistics of a random regular graph approach exactly the statistics allowed by the intended weak i.i.d. closure? This asks for more than tree-like uncolored neighborhoods.

The supplied assessment is explicitly provisional OPEN-TRIAGE: it reports no verified published result establishing the stated closure target, says its search cap was reached, and does not claim a definitive current-status review. It mentions the classical local weak comparison with the $d$-regular tree and describes local-global/measure-scaling convergence as a stronger notion capturing regularity. Its wording reverses the usual direction in the local-weak sentence; the conventional comparison is finite random graphs tending locally to the tree, not finite graphs being limits of a fixed tree. No local-global result is inferred from that comparison.
\subsection{Sources}
\begin{itemize}
\item UnsolvedMath permanent ID 1200019, source code catalogue entry, ``Some Questions --- Question 19.'' Registry curation and metadata: Ulam AI and UnsolvedMath Contributors, CC BY 4.0; updated 2026-08-06.
\item Abert, \emph{Some questions} (2010), Question 19: \url{https://www.renyi.hu/~abert/questions.pdf}. Supplied manual transcription accessed 2026-07-29. The author list was inspected according to the metadata, but current status was not independently verified there; status/release and rights review are NEEDS\_REVIEW.
\item The supplied provisional research assessment contains no additional authoritative reference for either the convergence target or an asserted resolution. Its unsuccessful capped search is reported as a limitation, not evidence proving openness.
\end{itemize}
\end{document}
